\chapter{Заключение}
\label{sec:zaver}

Цель этой бакалаврской работы — спроектировать и реализовать переносное устройство для тестирования и мониторинга компьютерных сетей, которое содержит генератор и анализатор трафика протоколов IPv4 и IPv6 и позволяет измерять такие параметры сети, как задержка, пропускная способность и потери. Эта цель достигнута.

В теоретической части описана проблематика генерирования сетевого трафика, включая разницу между измерением в битах и в кадрах в секунду, рекомендуемые размеры кадра и расчёт джиттера, и проблематика мониторинга сетей, включая фильтрацию трафика в ядре, сводные статистики и пределы самого захвата. Кроме того, разобраны адресация и автоконфигурация протокола IPv6 и технология пространств имён, на которой построен прибор.

В практической части создан прототип на основе одноплатного компьютера Raspberry Pi с сенсорным дисплеем, двумя внешними сетевыми адаптерами и аккумуляторным модулем. После включения прибор сразу запускает управляющее приложение и в типичном режиме работает от аккумулятора 6 часов 43 минуты. Он содержит анализатор трафика со сводными статистиками, разбором отдельных кадров, фильтрами в ядре и сохранением в формат pcap, генератор пакетов шести типов, измерение задержки тестом ping и измерение пропускной способности, потерь и джиттера инструментом iperf3, включая развёртку по размерам кадра и режим сервера. Для IPv6 он предлагает аудит сообщений Router Advertisement и поиск соседей с вычислением их адресов.

Решение проверено измерениями — и на перемычке между двумя портами прибора, и против чужой станции под Linux, и против ноутбука под Windows. Анализатор насчитал столько же кадров, сколько независимый пересчёт сохранённого файла, и его разбор совпал с анализатором Wireshark. Генератор доставил все отправленные пакеты во второе пространство имён. Пропускная способность по TCP против ноутбука составила 273--276~Мбит/с для обоих протоколов IP.

\section*{Собственный вклад автора}
\addcontentsline{toc}{section}{Собственный вклад автора}

Первый вклад — архитектура прибора, в которой оба тестовых интерфейса логически разделены пространствами имён ядра Linux, а каждая измерительная функция — отдельный скрипт, который работает в пространстве выбранного порта и передаёт результат управляющему приложению в машиночитаемом виде. Разделение устраняет конфликты маршрутизации между двумя интерфейсами в одном сегменте и проверено измерением: попытка обратиться к цели из неправильного пространства заканчивается недоступностью маршрута, а пакеты генератора попадают во второе пространство только по кабелю. Каждый результат сохраняется вместе со временем и контрольной суммой использованного скрипта, так что его происхождение можно проследить.

Второй вклад — разбор автоконфигурации адресов IPv6 на стороне анализатора. Показано, что dnsmasq при диапазоне, заданном одним адресом, опускает у объявленного префикса флаг A, и что после добавления ключевого слова \texttt{slaac} ядро станции само составляет адрес без какого-либо процесса в пространстве пользователя; контрольный порт с исходной настройкой подтвердил, что другая причина на результат не влияла. На этот разбор опирается модуль аудита, который разбирает сообщение Router Advertisement, оценивает его по требованиям стандартов и различает отправку собственного порта, второго порта того же прибора и чужого маршрутизатора, а также поиск соседей, который на одном и том же кабеле показал, что IPv6 находит другую сторону там, где IPv4 не находит ничего.

Третий вклад — пара генератора и анализатора трафика с проверенной точностью и измеренными пределами. Анализатор без потерь захватывает около 22\,500 кадров в секунду, генератор отправляет до 84\,000 кадров в секунду, и оба предела задаются производительностью процессора, а не объёмом данных. Измерения показали также, что на перемычке между портами пропускная способность вдвое ниже, потому что каждый кадр проходит шину USB дважды, и что адаптер прибора выбрасывает кадры UDP, если они приходят пачками быстрее, чем шина успевает их передать. Эти выводы определяют, как толковать измеренные числа.

Четвёртый вклад — воплощение устройства как самостоятельного прибора. Он работает от собственного аккумулятора с измерением заряда и штатным выключением при разряде, управляется пальцем на резистивном дисплее с элементами управления, рассчитанными на касание, и после подключения к чужой сети сам ничего не отправляет. Подключённой станции он выдаёт только адрес, а не шлюз по умолчанию, поэтому не нарушает её существующее подключение к Интернету. Для не прошедших тестов он объясняет оператору причину и советует, что делать дальше.

\section*{Ограничения и рекомендации по дальнейшему развитию}
\addcontentsline{toc}{section}{Ограничения и рекомендации по дальнейшему развитию}

Границы прибора нужно очертить честно. Ethernet использованной платформы подключён через шину USB 2.0, поэтому прибор достигает не более примерно 276~Мбит/с; он пригоден, чтобы проверить работоспособность соединения, сравнить протоколы в одинаковых условиях и выявить грубые неисправности, но не для сертификации гигабитной линии. Устранение этого ограничения потребовало бы перехода на более новое поколение одноплатного компьютера с контроллером Ethernet не на шине USB. Это же касается приёма пачечного трафика UDP и предела анализатора в кадрах в секунду.

Для дальнейшего развития есть несколько направлений. Ближайшее — корпус, в котором все компоненты будут жёстко закреплены; прибор уже работает от аккумулятора, но его части пока соединены свободно. Второе направление — автоматический тест розетки одной кнопкой, который после подключения к сети проверил бы линию, доступность сервера DHCP и маршрутизатора IPv6, шлюз по умолчанию и разрешение имён, как это делают коммерческие тестеры. Далее — определение свойств порта, то есть согласованной скорости, данных протоколов обнаружения соседних устройств и меток VLAN, затем определение наибольшего размера пакета, проходящего по пути, экспорт результатов на носитель и уточнение сенсорного управления калибровкой резистивной панели. Модуль аудита IPv6 подсказывает и направление безопасности: проверки, отработанные в этой работе на намеренно испорченном сообщении, в рабочей сети смогли бы обнаружить неправильно настроенный или несанкционированно подключённый маршрутизатор.
